\section{Model: Voting}

\subsection{Theory of Voting}
Voting classifiers are an ensemble learning technique that combine predictions from multiple base models to produce a final prediction. The approach is based on the principle that aggregating diverse models can reduce individual model biases and improve overall performance. Voting classifiers operate in two primary modes:
\begin{itemize}
    \item \textbf{Hard Voting}: Each base model casts a vote for a class label, and the final prediction is determined by majority voting.
    \item \textbf{Soft Voting}: Base models provide probability scores for each class, and the final prediction is based on the averaged probabilities across all models, typically selecting the class with the highest average probability.
\end{itemize}
Soft voting is often preferred when base models are well-calibrated, as it leverages confidence scores to make more informed decisions. Voting is effective when base models are diverse and complementary, as it mitigates weaknesses of individual models while preserving their strengths.

\subsection{Implementation}
We implemented two voting classifiers to enhance sentiment classification performance by combining predictions from multiple models. The voting process is managed by the \texttt{train\_voting\_classifier} function, which operates as follows:

\begin{itemize}
    \item \textbf{Training}: The function accepts a dictionary of base model classes, their hyperparameters, training features (\texttt{X}), labels (\texttt{y}), a feature extraction method, and a voting type (`soft' for probability averaging or `hard' for majority voting). It trains each base model independently, performs 5-fold cross-validation to evaluate performance (accuracy, ROC AUC, F1-score, precision, recall), and saves the trained voting classifier to a specified path. Performance and loss plots are generated.
    \item \textbf{Prediction}: The trained voting classifier predicts on test features by averaging probabilities (soft voting) or taking a majority vote (hard voting) from all base models to produce the final prediction.
\end{itemize}

The pseudocode for the \texttt{train\_voting\_classifier} function is as follows:

\begin{algorithm}[H]
\caption{\texttt{train\_voting\_classifier}}
\begin{algorithmic}[1]
\Procedure{TrainVotingClassifier}{model\_dict, param\_dict, feature\_method, X, y, voting\_type, model\_save\_path, img\_save\_path, img\_loss\_path}
    \If{file exists at model\_save\_path}
        \State \Return load model from model\_save\_path
    \EndIf
    \State base\_models $\gets$ empty list
    \For{each model\_name in model\_dict}
        \State model $\gets$ instantiate model\_dict[model\_name] with param\_dict[model\_name]
        \State append (model\_name, model) to base\_models
        \State handle exceptions by skipping invalid models
    \EndFor
    \If{length of base\_models < 2}
        \State \Return None
    \EndIf
    \State voting\_clf $\gets$ VotingClassifier(base\_models, voting\_type)
    \State Initialize lists for accuracy, roc\_auc, f1, precision, recall, train\_loss, val\_loss
    \State k\_fold $\gets$ KFold(n\_splits=5)
    \For{each train\_idx, val\_idx in k\_fold.split(X)}
        \State X\_train, X\_val $\gets$ X[train\_idx], X[val\_idx]
        \State y\_train, y\_val $\gets$ y[train\_idx], y[val\_idx]
        \State Fit voting\_clf on X\_train, y\_train
        \State Compute train\_loss and val\_loss using get\_training\_loss
        \State Append losses to train\_loss, val\_loss
        \State val\_preds $\gets$ voting\_clf.predict(X\_val)
        \State Append metrics (accuracy, roc\_auc, f1, precision, recall) for val\_preds, y\_val
    \EndFor
    \State Fit voting\_clf on full X, y
    \State Save voting\_clf to model\_save\_path
    \State Plot metrics and losses to img\_save\_path, img\_loss\_path
    \State \Return voting\_clf
\EndProcedure
\end{algorithmic}
\end{algorithm}

\subsection{Model 1: Voting Multiple Logistic Regression}
This model combines 10 Logistic Regression models, each with identical hyperparameters but different random seeds for diversity. The goal is to evaluate whether aggregating multiple strong Logistic Regression models (best-performing in prior experiments) improves results. Soft voting is used.

\textbf{Hyperparameters:}
\begin{verbatim}
{
    "model": "LogisticRegression",
    "instances": 10,
    "hyperparameters": {
        "penalty": "l2",
        "C": 0.1,
        "max_iter": 3000,
        "random_state": "range(0, 9)"
    }
}
\end{verbatim}

\textbf{Performance Metrics:}
\begin{table}[H]
\centering
\caption{Training and Testing Performance of Voting Multiple Logistic Regression}
\begin{tabular}{|l|c|c|}
\hline
\textbf{Metric} & \textbf{Training (Avg)} & \textbf{Testing} \\ \hline
Accuracy & 0.7440 & 0.7373 \\ \hline
ROC AUC & 0.7424 & 0.7356 \\ \hline
F1 Score & 0.7615 & 0.7568 \\ \hline
Precision & 0.7318 & 0.7232 \\ \hline
Recall & 0.7939 & 0.7938 \\ \hline
\end{tabular}
\end{table}

\begin{table}[H]
\centering
\caption{Testing Classification Report for Voting Multiple Logistic Regression}
\begin{tabular}{|l|c|c|c|}
\hline
\textbf{Class} & \textbf{Precision} & \textbf{Recall} & \textbf{F1-Score} \\ \hline
0.0 & 0.76 & 0.68 & 0.71 \\ \hline
1.0 & 0.72 & 0.79 & 0.76 \\ \hline
Weighted Avg & 0.74 & 0.74 & 0.74 \\ \hline
\end{tabular}
\end{table}

\subsection{Model 2: Voting Selected Models}
This model combines 4 diverse models that performed well in prior experiments: Logistic Regression, XGBoost, MLP, and GaussianNB, each configured with optimized hyperparameters. Soft voting is used to aggregate predictions, testing whether diverse models yield better performance.

\textbf{Hyperparameters:}
\begin{verbatim}
[
    {
        "model": "LogisticRegression",
        "hyperparameters": {
            "penalty": "l2",
            "C": 0.1,
            "max_iter": 1000
        }
    },
    {
        "model": "XGBoost",
        "hyperparameters": {
            "n_estimators": 150,
            "learning_rate": 0.1,
            "max_depth": 15
        }
    },
    {
        "model": "MLP",
        "hyperparameters": {
            "hidden_layer_sizes": [100],
            "activation": "logistic",
            "solver": "sgd",
            "alpha": 0.01,
            "batch_size": 32,
            "max_iter": 2000
        }
    },
    {
        "model": "GaussianNB",
        "hyperparameters": {
            "priors": [0.3, 0.7],
            "var_smoothing": 1e-9
        }
    }
]
\end{verbatim}

\textbf{Performance Metrics:}
\begin{table}[H]
\centering
\caption{Training and Testing Performance of Voting Selected Models}
\begin{tabular}{|l|c|c|}
\hline
\textbf{Metric} & \textbf{Training (Avg)} & \textbf{Testing} \\ \hline
Accuracy & 0.7284 & 0.7273 \\ \hline
ROC AUC & 0.7276 & 0.7263 \\ \hline
F1 Score & 0.7408 & 0.7415 \\ \hline
Precision & 0.7286 & 0.7244 \\ \hline
Recall & 0.7539 & 0.7594 \\ \hline
\end{tabular}
\end{table}

\begin{table}[H]
\centering
\caption{Testing Classification Report for Voting Selected Models}
\begin{tabular}{|l|c|c|c|}
\hline
\textbf{Class} & \textbf{Precision} & \textbf{Recall} & \textbf{F1-Score} \\ \hline
0.0 & 0.73 & 0.69 & 0.71 \\ \hline
1.0 & 0.72 & 0.76 & 0.74 \\ \hline
Weighted Avg & 0.73 & 0.73 & 0.73 \\ \hline
\end{tabular}
\end{table}

\subsection{Conclusion}
The voting model with multiple Logistic Regression models (\texttt{voting\_model\_multiple\_lr}) outperforms the voting model with selected diverse models (\texttt{voting\_model\_all\_models}) in both training and testing phases, achieving higher accuracy (0.7373 vs. 0.7273) and F1 score (0.7568 vs. 0.7415). This suggests that combining multiple instances of a strong model (Logistic Regression) is more effective than aggregating diverse models in this case.